% Tailored draft for ontinue-cloud-native-developer - base: powerapps-developer + full-stack
\documentclass[a4paper]{article}
    \usepackage{amsmath}
    \usepackage{amssymb}
    \usepackage{textcomp}
    \usepackage[utf8]{inputenc}
    \usepackage[T1]{fontenc}
    \textheight=10.8in
    \pagestyle{empty}
    \raggedright
    \usepackage{fontawesome}
    \usepackage{hyperref}
    \usepackage[left=0.5in, right=0.4in, top=0.4in, bottom=0.3in]{geometry}
    \linespread{0.92}
\def\bull{\vrule height 0.8ex width .7ex depth -.1ex }
\newcommand{\lineunder} {
    \vspace*{-8pt} \\
    \hspace*{-18pt} \hrulefill \\
}
\newcommand{\header} [1] {
    {\hspace*{-18pt}\vspace*{6pt} \textsc{#1}}
    \vspace*{-6pt} \lineunder
}
    \begin{document}
\vspace*{-40pt}
\vspace*{-10pt}
\begin{center}
	{\Huge \scshape {Muhammad Muizz Zafar}}\\
	\faSend \space \href{mailto:muizzafar11@gmail.com}{muizzafar11@gmail.com} $\cdot$ \faPhoneSquare \space 905-805-9142 $\cdot$ \faLinkedinSquare \space \href{https://www.linkedin.com/in/muizz-zafar/}{linkedin.com/in/muizz-zafar/}
	\\
\end{center}
\vspace{-5pt}
\begin{center}
\textbf{Cloud-Native Developer | .NET Core | Power Platform | REST APIs | Automation}
\end{center}
\begin{center}
Developer combining .NET Core APIs with Power Platform apps and Power Automate workflows, Docker and cloud pipelines.
\end{center}
\header{Skills}
\begin{tabular}{ l l }
	Programming Languages:        & C\#, ASP .NET, Python, Java, JavaScript, Kotlin, SQL, HTML/CSS, PHP, Flutter, C++\\
	Tools/Frameworks: & Power Platform, Power BI, AWS, Splunk, Flask, React, Docker, Git, Firebase, Linux \\
 \end{tabular}
\vspace{2mm}
\header{Education}
\textbf{University of Ottawa}\hfill Ottawa, ON\\
Bachelor of Applied Science, Software Engineering \textit{GPA: 3.7/4.0} \hfill Sept 2020 - Dec 2024\\
\vspace{2mm}
\header{Work experience}
\vspace{1.5mm}
\textbf{Software \& Database Analyst} \hfill May 2024 - Present\\
\textit{Elections Canada} \hfill Ottawa, ON\\
\vspace{-2mm}
\begin{itemize} \itemsep -2pt
\item Leveraged \textbf{Power Automate} to support and streamline \textbf{ETL validation workflows} by coordinating iterative testing processes, triggering \textbf{Python} comparison scripts, and automating pass/fail notifications with links to validation reports for stakeholders
\item Built internal \textbf{data analytics dashboards} using \textbf{Power BI} to display database health, operational metrics, validation trends, and data quality issues for directorate and senior management.
\item Implemented a reusable ETL validation framework using \textbf{Python} and Apache Spark, enabling automated data validation and improving data quality assurance across multiple environments.
\item Coordinated version-controlled \textbf{database deployments} and release processes, improving traceability and reliability of production data infrastructure across shared environments.
\item Collaborated in an \textbf{Agile cross-functional} team with stakeholders to improve ETL data workflows, release quality, and production reliability.
\end{itemize}
\textbf{Software Developer} \hfill September 2023 - April 2024\\
\textit{University of Ottawa} \hfill Ottawa, ON\\
\vspace{-2mm}
\begin{itemize} \itemsep -2pt
    \item Supported development and maintenance of \textbf{Microsoft Power Platform applications}, including \textbf{Power Automate}, to enhance workflow automation, data processing, and operational efficiency.
    \item Developed \textbf{Power Platform solutions} including \textbf{Canvas Apps and Model-Driven Apps} integrated with \textbf{C\# backend services} to support event management, user interactions, and data-driven workflows
    \item Implemented automated workflows using \textbf{Power Automate}, leveraging \textbf{HTTP triggers, scheduled flows, and API integrations} to enable real-time notifications, approvals, and data synchronization across systems
    \item Designed end-to-end integrations between \textbf{Power Apps and REST APIs}, enabling seamless communication between frontend applications and backend services
    \item Redesigned the University of Ottawa's emergency alert website with \textbf{C\# (.Net Core)} and \textbf{Angular}, introducing a modern, interactive UI. Used by over \textbf{20,000} students daily.
    \item Optimized existing \textbf{C\#} REST APIs by reducing \textbf{LINQ} query execution time for event data lookup by \textbf{30 \%} and lines of code by \textbf{40 \%}
    \item Participated in end-to-end \textbf{Software Development Lifecycle (SDLC)} processes, from requirements gathering and system design to implementation, testing, and production deployment
\end{itemize}
\textbf{Cloud Developer Intern} \hfill May 2023 - August 2023\\
\textit{Trend Micro Inc.} \hfill Ottawa, ON\\
\vspace{-2mm}
\begin{itemize} \itemsep -2pt
    \item Utilized \textbf{AWS Lambda} to monitor client files, images and code uploaded to the cloud and scan for vulnerabilities using Anchore APIs and Grype
        \item Programmed robust error handling for \textbf{AWS Lambda} functions in GoLang using goroutines to ensure proper timeout and recovery, reducing data corruption and improving service reliability
\end{itemize}
\textbf{Cloud Developer Intern} \hfill September 2022 - December 2022\\
\textit{Trend Micro Inc.} \hfill Ottawa, ON\\
\vspace{-2mm}
\begin{itemize} \itemsep -2pt
        \item Developed a serverless automation solution using \textbf{AWS Lambda, S3, and Athena crawlers} to trigger metadata updates on data changes, significantly reducing manual intervention and operational overhead
        \item Configured \textbf{Splunk} and \textbf{AWS CloudWatch} alarms for proactive monitoring of test logs and crashes, enhancing system reliability and reporting.
\end{itemize}
\textbf{Full Stack Developer Intern} \hfill January 2022 - April 2022\\
\textit{Applied Mind Inc.} \hfill Ottawa, ON\\
\vspace{-2mm}
\begin{itemize} \itemsep -2pt
\item Containerized a full-stack application with \textbf{Docker} for reliable \textbf{production deployments}, and automated event detection by adding events to \textbf{PostgreSQL} using \textbf{Python} and \textbf{GraphQl}, reducing testing time by \textbf{10 seconds}.
\end{itemize}
\vspace{5mm}
\end{document}
